\chapter{Listening to the place}\label{ch:acoustics}

The microphone records more than words. Every push-to-talk clip also says something about the space around
the phone, and a short chirp at the end of each report measures the distance to the walls. This chapter
describes those clues and how they are matched against the real buildings. It also says honestly how little
they locate on their own.

\section{Four acoustic clues}

\begin{table}[htbp]\centering\small
\caption{Acoustic clues and what they give.}\label{tab:aclues}
\begin{tabularx}{\linewidth}{>{\raggedright}p{3.2cm}>{\raggedright}p{4.6cm}X}
\toprule
Clue & Source & Gives\\
\midrule
Reverberation from the voice & Every push-to-talk clip & Decay time: $\sim$0.1 s outdoors, 0.3--0.6 s small room, 1--5 s hall or church\\
Chirp echo ranging & A 0.5 s sweep (800 Hz--12 kHz) played at the end of each report, or ``Listen 4 s'' & Impulse response: RT60, direct-to-reverberant ratio, the delay of every echo, so the distance to every wall ($d = 343\,\tau/2$)\\
Background level & The operator's own voice as loudness reference (about 76 dB(A) at the phone) & Noise floor in dB(A), $\pm$7 dB, e.g. 58 dB(A) for a residential street\\
Sound scene and sources & YAMNet + AST on the quiet part of the clip, or under the speech at half weight & Inside small / large / outside urban / rural; trains, road traffic, bells, machinery, birds\\
\bottomrule
\end{tabularx}
\end{table}

\nanofig{echo}{Chirp echo ranging in a street: the phone plays a sweep, and each wall returns it after a delay that gives its
distance. One microphone gives distance, never direction.}
{Top-down urban street illustration, 16:9, flat vector, plain white background, buildings as grey blocks with darker roofs.
A soldier with a phone in the middle of a narrow street. Concentric cyan rings radiating from the phone. Two building
facades highlighted, one at 6 m labelled ``echo 6.0 m'' and one across the street at 9 m labelled ``echo 9.0 m'', a
thin label ``street width 15 m (wall-to-wall echo)''. In the corner an inset graph titled ``echogram'' showing a large
first spike labelled ``direct'' and two smaller spikes labelled ``6 m'' and ``9 m'' on a distance axis in metres, with a
shaded band near 1 m labelled ``ground''. Exactly these labels. Colour palette (use only these): navy \#0D1110, cyan \#1AA7C4, amber \#E09A2B, red \#E5484D, green \#3FA877, grey \#8B9A94, plus white for text on dark shapes.}

\section{From audio to clues}

\paragraph{Sound classes.} Two classifiers run side by side on the laptop. YAMNet (Google, 4 MB, CPU) and the
Audio Spectrogram Transformer (AST, 87M parameters, GPU) both output the 521/527 AudioSet classes, grouped into
what matters: rail, traffic, bells, machinery, nature, and the danger groups. Their scores are averaged 0.4/0.6.
Speech windows are removed first; on a voice report, sounds are read under the speech at half weight, and only
what \emph{was} heard counts.

\paragraph{Kind of space.} Reverberation, echoes and scene scores become probabilities over four settings: indoors
small room, indoors large hall, outdoors between buildings, outdoors open ground.

\paragraph{Impulse response.} \code{spatial.py} deconvolves the recorded chirp with the known sweep, finds the direct
sound, and lists the reflections after it. A reflection between 0.8 and 1.8 m is usually the ground or a ceiling
and isn't counted as a wall. In a reverberant room the first wall echo hides in the tail, so it's marked low
confidence and not used.

\section{Matching against the real place}

\code{surround.py} computes what the phone \emph{should} hear at every spot: 9 spots per 25 m cell, 245,520 in
all, cached in about 25 s.
\begin{itemize}
  \item \textbf{Expected echoes.} 36 rays per spot over a 1 m raster of the IGN footprints. A flat wall reflects
        the chirp from the foot of the perpendicular, so the local minima of distance along the rays are the expected
        echo distances. Two facing walls also give a wall-to-wall echo at the street width. Each measured echo is
        matched to the nearest expected one ($\sigma$ = 0.5 m + 5 \%). A false-alarm floor covers cars, trees and
        people. A big wall within 15 m that returned nothing costs a penalty.
  \item \textbf{Expected noise level.} Every OSM road and rail point is a source: 78 dB(A) at 10 m for a motorway, down
        to 65 for a tertiary road. Each falls off as a line source, loses 15 dB behind a building and 25 dB through a
        facade indoors. In the woods there is no loss for the first 20 m of trees, then 0.05 dB/m, capped at 15 dB.
  \item \textbf{Bits per clue.} The dashboard shows how much each clue narrowed the map, and the area still holding 90 \%
        of the probability with audio alone (the ``Audio match'' layer).
\end{itemize}

\section{Results}

\begin{table}[htbp]\centering\small
\caption{Sound classifiers on 85 labelled real recordings (ESC-50), threshold 0.15.}\label{tab:yamnetast}
\begin{tabular}{lrrrrr}
\toprule
Model & Recall & Precision & F1 & Top group right & False location alarms\\
\midrule
YAMNet & 92 \% & 73 \% & 0.81 & 46/50 & 4/35\\
AST & 80 \% & 98 \% & 0.88 & 45/50 & 0/35\\
\textbf{YAMNet + AST} & 92 \% & 87 \% & \textbf{0.89} & \textbf{48/50} & 2/35\\
\bottomrule
\end{tabular}
\end{table}

\begin{itemize}
  \item \textbf{Blind reverberation} from real sounds in rooms of known RT60: 0.6 s $\rightarrow$ 0.59 s,
        1.2 s $\rightarrow$ 1.12 s, 2.4 s $\rightarrow$ 2.14 s. Small rooms are less precise (0.3 s $\rightarrow$ 0.38 s).
  \item \textbf{Echoes} come back within 0.3 m of the true walls on clips simulated from the real buildings; the noise level
        is within about 3 dB of the model.
  \item \textbf{Ask-back loop} (simulated): median error 102 m with landmarks only, 83 m with the space and sound questions,
        62 m with a Listen clip. That last number is an upper bound, since the clip is drawn from our own model.
  \item \textbf{With 3 landmarks and the sun}, one voice clip with its end chirp raises ``best guess within 75 m'' from 12 \%
        to 18 \% and cuts the median error of the best guess from 1,060 m to 810 m.
\end{itemize}

\begin{caveat}
Audio alone narrows the \emph{kind} of spot, not the place. The true spot ranks in the top 6 \% of the map (median) and
audio removes about 0.8 bits. One microphone hears distance, not direction, and echoes repeat all along a uniform
street, so inside a known 300 m area audio alone is only about 1.5$\times$ better than chance. Scores and thresholds
are not calibrated on real recordings yet. Playing a chirp makes noise; it can be turned off, and reverberation from
the voice still works passively.
\end{caveat}

\section{Heart rate from a chest recording}

\code{heart.py} measures the pulse when the phone's microphone is pressed on the casualty's chest for 10 s. It takes the
envelope of the heart-sound band (25--200 Hz) and its autocorrelation, and checks every multiple of each candidate period
so the ``lub-dub'' gap isn't taken for a beat. On synthetic chest recordings, 96 \% of clean ones were within 5 bpm, and
readings flagged reliable were within 5 bpm 92 \% of the time. Noisy ones are marked ``unclear: press firmly, repeat''.
A voice recording does not carry a usable heart rate, so the dashboard only shows a value from a chest recording, with its
age. Bands: slow $<$ 50, normal, fast $>$ 100, shock range $>$ 120, severe $>$ 140 bpm (placeholders for the mentor).
